\documentclass[11pt]{article}

\usepackage{series}
\usepackage{amsmath,amssymb,amsthm}
\usepackage{geometry}
\usepackage{hyperref}
\usepackage{booktabs}

\geometry{margin=1in}

% --- Metadata ---
\title{Theta Closure in Modal Triplet Theory III:
Conditional Twistor--Action Matching and Normalization Audit}
\author{Peter Nero}
\date{September 2026 (Version 3)}

\begin{document}
\maketitle

\begin{abstract}
We audit a conditional Route~B representation of the leading--order nonabelian
overlap integrals $I_2^{(0)}$ and $I_3^{(0)}$ in Modal Triplet Theory (MTT)
using the self--dual Yang--Mills (SDYM) twistor corner. A standard SDYM
twistor action is holomorphic BF-type; it does not by itself provide the
positive Hermitian $L^2$ norm or relative coupling normalization used by
Route~A. We therefore isolate a reconstruction-and-norm bridge as an explicit
assumption. The $SU(2)$ comparison additionally uses
$dA_{\mathrm{dir}}=2\omega_{\mathrm{FS}}$ and the effective round-$S^2$ lens
model, while the $SU(3)$ comparison uses the declared auxiliary nilmanifold,
color harmonic, and factorization assumption of Paper~II. Under these shared
inputs Route~B reproduces
$I_2^{(0)}=4\pi(f_2R_{\mathrm{lens}})^2$ and
$I_3^{(0)}=\int\|\chi_{\mathrm{col}}\|^2d\mu=c$. The result is therefore a
conditional dictionary check, not an independent derivation or selection of
the overlap values and not a completed nonabelian $\Theta$--closure theorem.
\end{abstract}

\section*{Version 3 Revision Note}
\label{sec:revision-v3}
\begin{description}
\item[Supersedes.] Version~2 of this paper; released identifiers are retained.
\item[Reason.] The calibration boundary was present, but the selected-source
owner results and their distinct scopes were not explicitly connected to it.
\item[Resolution.] Version~3 imports the exact Gate~1, three-cycle and
large-gauge-kernel conclusions through their owners, and also restores period normalization in the final color comparison.
\item[Retained.] The existing calibrated or conditional result, its numerical
inputs and all earlier revision notes are retained.
\item[Open boundary.] A selected physical source, its action normalization and
the paper-specific execution inputs remain necessary for a held-out prediction.
\end{description}


\section*{Revision note for this edition}
\addcontentsline{toc}{section}{Revision note for this edition}
\begin{description}
\item[Supersedes.] \emph{Theta Closure in Modal Triplet Theory III: Twistor
Action Matching and Independent Normalization}, first edition.
\item[Reason.] Route~B was described as an independent normalization although
its area convention, effective lens model, nil harmonic, and factorization
were inherited from Route~A. In addition, the displayed twistor action was a
schematic Hermitian curvature norm rather than a standard holomorphic
BF-type SDYM action.
\item[Resolution.] Version~2 exposes every shared input and proves only the
conditional correspondence obtained after a reconstruction-and-norm bridge is
imposed. The standard twistor action and the additional Hermitian norm bridge
are now kept distinct.
\item[Retained result.] Under those declared bridge conventions, the twistor
and direct descriptions use compatible leading overlap functionals.
\item[Remaining boundary.] An independently selected normalization or q79
source theorem is still required for redundant closure evidence.
\end{description}


\section{Relation to Papers I and II}

This paper is the third component of the conservative $\Theta$--closure program:
\begin{itemize}
\item \textbf{Paper I} defines the selected common-scheme gauge profile and
the corresponding overlap-ratio targets.
\item \textbf{Paper II} realizes those targets in a calibrated effective
round-$S^2$/nilmanifold ansatz (Route~A).
\item \textbf{Paper III} tests whether the same quadratic overlap functional
is reproduced by the twistor--encoded gauge action (Route~B), and audits which
normalization inputs remain shared with Route~A.
\end{itemize}

Agreement between Routes~A and~B is a conditional internal consistency check.
It is not statistically or logically independent evidence when the routes
share the internal metric, harmonic, normalization bridge, or target profile.

\section{Central picture and reading guide}

The central picture is a dictionary between two descriptions. Twistor theory
encodes a self-dual gauge field by holomorphic bundle data on twistor space.
Route~A encodes a gauge coefficient by a weighted norm of an internal
representative. These are not automatically the same mathematical object. To
compare them one must provide a bridge that turns reconstructed twistor data
into the Hermitian quadratic form used by the effective gauge action.

The paper therefore has three logical stages. First, it states the standard
Penrose--Ward and holomorphic-action input. Second, it declares the extra norm,
coupling, direction-sphere, and color-factorization bridges. Third, it checks
that those bridges reproduce the two Route~A formulas. The final equality is
useful as a compatibility audit, but it cannot be counted as a second
measurement or a second source of the numerical values.

\section{Introduction}

Modal Triplet Theory (MTT) proposes that observable four--dimensional physics
arises from a coherent projection of a higher--dimensional modal configuration
space. In this framework, inverse squared gauge couplings are represented by
overlap integrals of internal harmonic representatives. Whether the values of
those integrals are selected by MTT geometry or calibrated to an observed
profile is a separate theorem question.

Paper~I transports the measured Standard Model gauge profile in a common
scheme to $Q=M_t=172.5590883453979~\mathrm{GeV}$. Paper~II maps the resulting
ratios into an auxiliary Lens--Nil operator model: a round-$S^2$ effective
lens base and a compact Heisenberg nilmanifold, with one shared $U(1)$
phase/holonomy datum reused across the lanes. It does not use the
seven-dimensional product
$S^1_{\mathrm{cen}}\times L(3,1)\times
\Gamma\backslash\mathrm{Nil}_3$. Its parameters constitute a calibrated
ansatz-level realization, not a unique geometry selection or an identification
with the selected q79/Fu--Yau compactification.

The purpose of this paper is to determine precisely what the twistor--action
construction (Route~B) adds. It represents the quadratic gauge norm through
the declared reconstruction-and-norm bridge, not the SDYM action alone.
It does not, by itself, select
the map from the Fubini--Study form to the effective lens area, the internal
color fiber, or the color harmonic. Those bridges are stated as assumptions
rather than hidden inside a claim of independent normalization.

We show that:
\begin{itemize}
\item the declared fiber-reduction bridge assigns a quadratic $L^2$ norm
to the massless representative;
\item the declared direction-sphere bridge recovers the $4\pi$ coefficient
used by the effective $SU(2)$ lens-base model;
\item the selected color factorization reproduces the nilmanifold $SU(3)$
overlap functional; and
\item numerical equality with Route~A is conditional on these shared inputs,
while a quantitative bound on omitted coherence corrections remains open.
\end{itemize}

The result is a representation-level compatibility theorem, not an
independent value source or a completed nonabelian $\Theta$--closure theorem.

\section{Twistor corner and SDYM action}

\subsection{High--coherence twistor regime}

The twistor formulation of MTT identifies a high--coherence regime in which the
effective dynamics of the massless gauge sector reduces to self--dual Yang--Mills
(SDYM) theory encoded holomorphically on twistor space.

\begin{assumption}[Twistor--corner regime]
We assume the matching scale $\mu_\Theta$ lies in the high--coherence regime
for the massless gauge sector, so that:
\begin{enumerate}
\item SDYM provides the leading--order gauge dynamics;
\item massive corrections are suppressed by $O(\lambda_Q^{-1})$;
\item a specified reconstruction map identifies the twistor fields with the
spacetime self-dual gauge sector.
\end{enumerate}
\end{assumption}

\subsection{Twistor action for SDYM}

Standard twistor formulations of SDYM are holomorphic BF-type (or closely
related holomorphic Chern--Simons-type) actions
\cite{Mason2005,Movshev2008}. Schematically, on a suitable twistor region one
may write
\[
S_{\mathrm{SDYM}}[{\cal A},{\cal B}]
=
\frac{1}{g_{\mathrm{tw}}^2}
\int_{\mathbb{PT}}
\Omega\wedge\mathrm{Tr}\!\left({\cal B}\wedge{\cal F}^{(0,2)}\right),
\]
where the bundle weights of $\Omega$ and ${\cal B}$ are chosen so that the
integrand is globally defined. Here:
\begin{itemize}
\item $\mathbb{PT}$ is the relevant region of projective twistor space,
\item ${\cal F}^{(0,2)}$ is the antiholomorphic curvature of the $(0,1)$
connection ${\cal A}$,
\item ${\cal B}$ is the multiplier field imposing
${\cal F}^{(0,2)}=0$, and
\item $g_{\mathrm{tw}}$ is the twistor--action coupling.
\end{itemize}

The Penrose--Ward correspondence relates the resulting holomorphic bundle data
to the self-dual Yang--Mills equations. It does not, without further choices,
produce a positive Hermitian fiber norm or fix its normalization relative to
the full four-dimensional Yang--Mills kinetic term. That additional datum is
the reconstruction-and-norm bridge introduced next.

\section{Reconstruction of the four--dimensional gauge action}

\subsection{Ward correspondence and spacetime fields}

By the Ward correspondence, holomorphic bundles satisfying the twistor-line
triviality and relevant reality conditions reconstruct SDYM fields on the
stated spacetime domain. A kinetic coefficient from fiber integration is
the additional assumption below, not a consequence of reconstruction alone.

\begin{assumption}[Reconstruction-and-norm bridge]
For the selected massless sector, assume that the reconstructed field is
assigned the quadratic four-dimensional functional
\[
Q_4[A] = \frac{1}{4g_{\mathrm{eff}}^2}
\int_{Y^4} F_{\mu\nu}F^{\mu\nu}\,d\mathrm{vol}_4.
\]
Assume further that its coefficient is related to a chosen Hermitian fiber
metric by
\[
\frac{1}{g_{\mathrm{eff}}^2}
=
\frac{1}{g_{\mathrm{tw}}^2}
\int_{\mathcal F} \langle \psi, \psi\rangle\,d\mu_{\mathcal F},
\]
where:
\begin{itemize}
\item $\mathcal F$ is the internal fiber associated with the gauge sector,
\item $\psi$ is the period/bridge-normalized twistor harmonic corresponding to
the massless mode, with any gauge-kinetic weight included in the displayed
measure.
\end{itemize}
\end{assumption}

\subsection{Matching of couplings}

Under this bridge assumption, the displayed fiber norm is the proposed
twistor analogue of the internal overlap integral $I_a^{(0)}$ defined in
Papers~I and~II. Neither the holomorphic SDYM field equation nor the
Fubini--Study area convention alone proves this coupling identity.

\section{Identification of
\texorpdfstring{$I_2^{(0)}$ and $I_3^{(0)}$}{I2(0) and I3(0)}}

\subsection{General matching principle}

Comparing the twistor reconstruction with the overlap formulation
\[
\frac{1}{g_a^2} = \frac{1}{g_{10}^2} I_a,
\]
we identify
\[
I_a^{(0)} \;\equiv\;
\int_{\mathcal F_a}
w_a^{(0)}\langle \psi_a^{(0)}, \psi_a^{(0)}\rangle\,
d\mu_{\mathcal F_a},
\]
where $\psi_a^{(0)}$ is the massless twistor harmonic for the $SU(a)$ sector.

\subsection{\texorpdfstring{$SU(2)$}{SU(2)} sector:
conditional Route~B normalization}

We evaluate the Route~B quadratic norm for the $SU(2)$ sector. This avoids the
internal spectral estimate used in Route~A, but the numerical coefficient still
depends on the coupling and direction-sphere bridge conventions stated below.

\subsubsection{Twistor fiber normalization}

In the high--coherence twistor corner, each spacetime point $x\in Y^4$
corresponds to a twistor line $L_x\cong\mathbb{CP}^1$.
We equip $L_x$ with the standard Fubini--Study K\"ahler form
$\omega_{\mathrm{FS}}$ normalized by
\begin{equation}
\int_{L_x}\omega_{\mathrm{FS}} = 2\pi.
\label{eq:FS_norm}
\end{equation}
This convention fixes the fiber measure once chosen. It does not by itself
fix the relative normalization between $g_{\mathrm{tw}}$, $g_{10}$, and the
internal MTT overlap.

\subsubsection{Twistor SDYM action}

We use the holomorphic BF-type SDYM action schematically as
\begin{equation}
S_{\mathrm{SDYM}}[{\cal A},{\cal B}]
=
\frac{1}{g_{\mathrm{tw}}^2}
\int_{\mathbb{PT}}
\Omega\wedge\mathrm{Tr}\!\left(
{\cal B}\wedge{\cal F}^{(0,2)}
\right),
\label{eq:twistor_action_SU2}
\end{equation}
with the required line-bundle weights understood. This action supplies the
holomorphic SDYM equations. The normalization of a positive Hermitian kinetic
functional is additional data and is not fixed merely by choosing
\eqref{eq:FS_norm}.

\subsubsection{Reconstruction to spacetime}

By the Ward correspondence, solutions of the twistor field equations correspond
to self--dual Yang--Mills fields on spacetime.
Under the reconstruction-and-norm bridge, restricting to the selected
massless sector assigns the four--dimensional quadratic functional
\begin{equation}
S_{4}[A]
=
\frac{1}{4g_{\mathrm{eff}}^2}
\int_{Y^4} F_{\mu\nu}F^{\mu\nu}\,d\mathrm{vol}_4,
\label{eq:4d_SU2}
\end{equation}
where the effective coupling is assumed to obey the fiber identity
\begin{equation}
\frac{1}{g_{\mathrm{eff}}^2}
=
\frac{1}{g_{\mathrm{tw}}^2}
\int_{L_x}
\|\psi^{(0)}\|^2\,d\mu_{L_x}.
\label{eq:SU2_fiber_norm}
\end{equation}
Here $\psi^{(0)}$ is the massless twistor harmonic corresponding to the
$SU(2)$ gauge zero mode.

\subsubsection{Evaluation of the fiber integral}

For this comparison we choose the $SU(2)$ massless representative to be
constant along the fiber and use the Hermitian normalization
\[
\|\psi^{(0)}\|_{L^2(L_x)}^2
=
\int_{L_x}\|\psi^{(0)}\|^2\,d\mu_{L_x}
=
\int_{L_x}\omega_{\mathrm{FS}}
=
2\pi,
\]
by \eqref{eq:FS_norm}.
Substituting into \eqref{eq:SU2_fiber_norm} gives
\begin{equation}
\frac{1}{g_{\mathrm{eff}}^2}
=
\frac{2\pi}{g_{\mathrm{tw}}^2}.
\label{eq:SU2_eff_coupling}
\end{equation}

\subsubsection{Identification with the MTT overlap}

In the MTT overlap formulation at the matching scale,
\[
\frac{1}{g_2^2}
=
\frac{1}{g_{10}^2}
\int_{B_2}w_2\langle\omega_2,\omega_2\rangle\,d\mu_{B_2}
=
\frac{1}{g_{10}^2}\,I_2^{(0)}.
\]
Comparing with \eqref{eq:SU2_eff_coupling}, we identify the leading--order
$SU(2)$ overlap as
\begin{equation}
I_2^{(0)}
=
\int_{B_2}w_2^{(0)}
\langle\omega_2^{(0)},\omega_2^{(0)}\rangle\,d\mu_{B_2}
=
4\pi (f_2R_{\mathrm{lens}})^2,
\label{eq:SU2_overlap_final}
\end{equation}
where the factor $4\pi$ arises from the area of the unit two--sphere of
directions associated with the twistor line.

\begin{theorem}[Conditional Route~B $SU(2)$ normalization]
Assume the SDYM fiber reduction, the coupling identification, and the
direction-sphere bridge $dA_{\mathrm{dir}}=2\omega_{\mathrm{FS}}$. For the
$SU(2)$ gauge sector the resulting leading--order overlap is
\[
I_2^{(0)} = 4\pi (f_2R_{\mathrm{lens}})^2,
\]
without using the internal Laplacian bound. The metric scale and the
direction-sphere bridge are nevertheless shared with the Route~A ansatz.
\end{theorem}

\begin{remark}
This gives a compatible representation of the quadratic norm after the norm
bridge has been imposed, but not an independent derivation or numerical value
source. The Hermitian fiber metric, coupling identification, direction-sphere
bridge, and effective lens metric remain declared inputs.
\end{remark}

\subsection{\texorpdfstring{$SU(3)$}{SU(3)} sector:
conditional color--twistor factorization}

We next incorporate the selected internal color fiber into Route~B. The
resulting identity is conditional because neither the fiber nor its massless
harmonic is selected by the twistor action in this paper.

\subsubsection{Color--twistor factorization of the massless mode}

In the auxiliary Route~A ansatz, the effective color support is the compact
Heisenberg nilmanifold $\Gamma\backslash\mathrm{Nil}_3$. The shared $U(1)$
circle is common phase/holonomy data reused across the gauge lanes; it is not
inserted here as a second Cartesian factor of the color support.
In the high--coherence regime, the massless $SU(3)$ gauge mode admits a
factorized representation
\begin{equation}
\Psi^{(0)}(Z,u)
=
\psi_{\mathrm{tw}}(Z)\otimes \chi_{\mathrm{col}}(u),
\label{eq:color_twistor_factorization}
\end{equation}
where:
\begin{itemize}
\item $Z\in L_x\simeq\mathbb{CP}^1$ is the spacetime twistor coordinate,
\item $u\in \Gamma\backslash\mathrm{Nil}_3$ is the internal color coordinate,
\item $\psi_{\mathrm{tw}}$ is the normalized massless twistor harmonic,
\item $\chi_{\mathrm{col}}$ is a normalized massless color harmonic on
$\Gamma\backslash\mathrm{Nil}_3$.
\end{itemize}

This factorization is an explicit Route~B assumption. The present paper does
not derive it, select the compact nilmanifold, or prove uniqueness of the color
harmonic from the twistor action alone.

\subsubsection{Normalization of the factorized harmonics}

The twistor harmonic $\psi_{\mathrm{tw}}$ is normalized using the standard
Fubini--Study measure $\omega_{\mathrm{FS}}$ on $\mathbb{CP}^1$:
\begin{equation}
\int_{L_x} \|\psi_{\mathrm{tw}}\|^2\, d\mu_{\mathrm{FS}}
=
\int_{L_x} \omega_{\mathrm{FS}}
=
2\pi.
\label{eq:SU3_twistor_norm}
\end{equation}

The color harmonic $\chi_{\mathrm{col}}$ uses the period normalization of
Paper~II. Its $L^2$ norm is then an output, not an imposed unit value:
\begin{equation}
\|\chi_{\mathrm{col}}\|_{L^2(\Gamma\backslash\mathrm{Nil}_3)}^2
:=
\int_{\Gamma\backslash\mathrm{Nil}_3}
\|\chi_{\mathrm{col}}\|^2\, d\mu_{\mathrm{nil}}.
\label{eq:SU3_color_norm}
\end{equation}

\subsubsection{Reconstructed quadratic normalization for
\texorpdfstring{$SU(3)$}{SU(3)}}

Substituting the factorized form \eqref{eq:color_twistor_factorization} into the
declared reconstruction-and-norm bridge gives the four--dimensional quadratic
functional
\begin{equation}
S_4[A]
=
\frac{1}{4g_{\mathrm{eff}}^2}
\int_{Y^4} F_{\mu\nu}F^{\mu\nu}\, d\mathrm{vol}_4,
\end{equation}
with effective coupling
\begin{equation}
\frac{1}{g_{\mathrm{eff}}^2}
=
\frac{1}{g_{\mathrm{tw}}^2}
\left(
\int_{L_x}\|\psi_{\mathrm{tw}}\|^2\, d\mu_{\mathrm{FS}}
\right)
\left(
\int_{\Gamma\backslash\mathrm{Nil}_3}
\|\chi_{\mathrm{col}}\|^2\, d\mu_{\mathrm{nil}}
\right).
\label{eq:SU3_eff_coupling}
\end{equation}

Using \eqref{eq:SU3_twistor_norm} and \eqref{eq:SU3_color_norm}, we obtain
\begin{equation}
\frac{1}{g_{\mathrm{eff}}^2}
=
\frac{2\pi}{g_{\mathrm{tw}}^2}
\int_{\Gamma\backslash\mathrm{Nil}_3}
\|\chi_{\mathrm{col}}\|^2\, d\mu_{\mathrm{nil}}.
\label{eq:SU3_eff_coupling_final}
\end{equation}

\subsubsection{Identification with the MTT overlap formulation}

In the MTT overlap formulation at the matching scale,
\begin{equation}
\frac{1}{g_3^2}
=
\frac{1}{g_{10}^2}
\int_{B_3}
w_3\langle \omega_3, \omega_3 \rangle\, d\mu_{B_3}.
\end{equation}

Comparing with \eqref{eq:SU3_eff_coupling_final}, we define the leading--order
$SU(3)$ overlap intrinsically by
\begin{equation}
I_3^{(0)}
:=
\int_{\Gamma\backslash\mathrm{Nil}_3}
\|\chi_{\mathrm{col}}\|^2\, d\mu_{\mathrm{nil}}.
\label{eq:SU3_overlap_def}
\end{equation}

For the isotropic left--invariant metric $a=b$ on
$\Gamma\backslash\mathrm{Nil}_3$ and choosing $\chi_{\mathrm{col}}$ to be a
canonical left--invariant harmonic one--form, this evaluates explicitly to
\begin{equation}
I_3^{(0)} = c.
\label{eq:SU3_overlap_value}
\end{equation}

\begin{theorem}[Conditional Route~B $SU(3)$ overlap identity]
In the high--coherence twistor regime, assuming the reconstruction-and-norm
bridge and canonical factorization of the
massless $SU(3)$ gauge mode into a spacetime twistor harmonic and an internal
color harmonic on $\Gamma\backslash\mathrm{Nil}_3$, the twistor--action
reduction yields the conditional leading--order identity
\[
I_3^{(0)}
=
\int_{\Gamma\backslash\mathrm{Nil}_3}
\|\chi_{\mathrm{col}}\|^2\, d\mu_{\mathrm{nil}},
\]
which evaluates to $I_3^{(0)}=c$ for the isotropic metric $a=b$ and the
Paper~II harmonic convention.
\end{theorem}

\begin{remark}
Together with the $SU(2)$ analysis, this supplies a conditional Route~B
compatibility check for the massless nonabelian sector. The assumed asymptotic
order $O(\lambda_Q^{-1})$ is not yet accompanied here by a coefficient or
uniform remainder bound and must not be read as a numerical error certificate.
\end{remark}


\subsection{Nil sector (\texorpdfstring{$SU(3)$}{SU(3)})}

For the $SU(3)$ sector, the twistor harmonic reconstructs to a left--invariant
massless gauge mode on the compact nilmanifold
$\Gamma\backslash\mathrm{Nil}_3$.

The twistor fiber integral reduces to the $L^2$ norm of the reconstructed
spacetime harmonic:
\[
I_3^{(0)} = \int_{\Gamma\backslash\mathrm{Nil}_3}
|\omega^{(0)}|^2\,d\mathrm{vol}.
\]

For the isotropic horizontal metric choice $a=b$, this evaluates to
\[
I_3^{(0)} = c,
\]
in agreement with the direct geometric computation of Paper~II.

\section{Consistency with the selected gauge profile}

Paper~I gives the selected common-scheme targets at
$Q=M_t=172.5590883453979~\mathrm{GeV}$:
\[
\frac{I_2}{I_1}=0.5110273\pm0.0001231,
\qquad
\frac{I_3}{I_1}=0.158335\pm0.001098,
\qquad
I_1=2\pi R_1.
\]
With the leading identities used in Papers~II and~III, these imply
\[
(f_2R_{\mathrm{lens}})^2=0.2555137\,R_1,
\qquad
c=0.9948493\,R_1.
\]
The former $5~\mathrm{TeV}$ profile and the values $0.560$, $0.229$,
$0.280R_1$, and $1.439R_1$ are withdrawn.

\begin{theorem}[Conditional Route~A/Route~B agreement]
Assume the common coupling-overlap convention, the selected Paper~I profile,
the Paper~II effective internal ansatz, the SDYM fiber reduction, the
direction-sphere normalization bridge, and color--twistor factorization. Then
both routes evaluate the same leading quadratic functional and give
\[
I_2^{(0)}=4\pi(f_2R_{\mathrm{lens}})^2,
\qquad
I_3^{(0)}=c.
\]
Consequently both reproduce the displayed profile ratios.
\end{theorem}

\begin{proof}
Under the listed assumptions, both constructions evaluate the $L^2$ norm of
the same selected massless representative. Equality follows from the shared
normalization bridges. It is an exact round-trip identity within the selected
ansatz, not an independent derivation of the input profile or internal metric.
\end{proof}

\subsection{Worked example: what the agreement does and does not test}

Set $R_1=1$ and import the calibrated Paper~II scales
\[
(f_2R_{\mathrm{lens}})^2=0.2555137,
\qquad
c=0.9948493.
\]
Route~A then gives
\[
I_2^{(0)}=4\pi(0.2555137)\approx3.210880,
\qquad
I_3^{(0)}=0.9948493.
\]
Route~B returns the same numbers because the direction-sphere bridge, the
nilmanifold harmonic, and the reconstruction-and-norm bridge identify its
fiber expressions with those same two functionals. The example checks that the
dictionary is internally consistent. It does not supply new data with which to
test the values $3.210880$ or $0.9948493$.

\section{Dependency audit}

The Route~B comparison shares the following inputs with Papers~I and~II:
\begin{enumerate}
\item the measured common-scheme gauge profile and $I_1=2\pi R_1$ convention;
\item the relation $g_a^{-2}=g_{10}^{-2}I_a$ and its common normalization;
\item the effective round-$S^2$ lens-base metric and scale
$f_2R_{\mathrm{lens}}$;
\item the compact nilmanifold, its metric parameter $c$, and the chosen color
harmonic; and
\item the identification of the twistor fiber norm with the selected internal
overlap, including $dA_{\mathrm{dir}}=2\omega_{\mathrm{FS}}$.
\end{enumerate}
Route~B independently supplies the holomorphic SDYM representation of the
self-dual field data. Its interpretation as the displayed quadratic norm is
conditional on the reconstruction-and-norm bridge. A genuinely independent value
derivation would have to select the bridges and internal representatives from
twistor-corner MTT data without importing the Route~A realization.

\section{Scope and limitations}

This paper establishes conditional leading-order compatibility between two
representations of the nonabelian overlap functional. It does not establish:
\begin{itemize}
\item a source theorem selecting the gauge profile or internal geometry;
\item an independent $SU(2)$ or $SU(3)$ numerical normalization;
\item a coefficient-level bound for corrections denoted
$O(\lambda_Q^{-1})$;
\item Yukawa/flavor structure, ultraviolet completion, or string embedding.
\end{itemize}

\section{Conclusion}

The SDYM twistor description and the direct internal calculation can be placed
on the same quadratic $L^2$ footing only after a reconstruction-and-norm
bridge is declared. Once that bridge and the shared geometry conventions are
imposed, the two descriptions agree exactly at leading order and reproduce the
selected Paper~I overlap profile. This checks the internal consistency of the
proposed dictionary; it is not a derivation of that dictionary from the
holomorphic action.

It is not a second independent determination of the numerical overlaps.
The $SU(2)$ coefficient uses the direction-sphere bridge and effective lens
metric; the $SU(3)$ identity uses the declared auxiliary nilmanifold and color
harmonic.
The strongest justified conclusion is therefore conditional
representation-level compatibility. Independent geometric selection and a
quantitative coherence-remainder estimate remain separate theorem targets.

% ===========================
\appendix
\section{Normalization-bridge audit: twistor data to overlap coefficient}
% ===========================

This appendix records the Route~B normalization dictionary without promoting
it to a theorem of the holomorphic SDYM action. The standard twistor
correspondence supplies the self-dual field data. A Hermitian norm, its
relation to the four-dimensional kinetic coefficient, and its identification
with the MTT internal overlap are additional bridge choices.

\subsection{Twistor geometry and canonical fiber measure}

Let $Y^4$ be a (local) conformally flat spacetime patch so that its twistor space
$\mathbb{PT}$ is a $\mathbb{CP}^1$ bundle over $Y^4$.
Write the projection as
\[
\pi:\mathbb{PT}\to Y^4.
\]
Each point $x\in Y^4$ corresponds to a twistor line
\[
L_x \cong \mathbb{CP}^1.
\]

We use the standard Fubini--Study K\"ahler form $\omega_{\mathrm{FS}}$ on $\mathbb{CP}^1$
normalized by
\begin{equation}
\int_{\mathbb{CP}^1}\omega_{\mathrm{FS}} = 2\pi.
\label{eq:FSnorm}
\end{equation}
This fixes the fiber measure within the adopted Fubini--Study convention.

\begin{remark}
Equation \eqref{eq:FSnorm} is a normalization convention: any rescaling would
correspond to a rescaling of the twistor coupling, and therefore must be fixed once.
We fix it as above; a separate bridge is still required to compare the twistor
coupling and fiber norm with the MTT internal overlap.
\end{remark}

\subsection{Holomorphic SDYM action}

A standard twistor-space action for self-dual Yang--Mills has the schematic
holomorphic BF form
\begin{equation}
S_{\mathrm{SDYM}}[{\cal A},{\cal B}]
=
\frac{1}{g_{\mathrm{tw}}^2}
\int_{\mathbb{PT}}
\Omega\wedge\mathrm{Tr}\!\left(
{\cal B}\wedge{\cal F}^{(0,2)}
\right),
\label{eq:twistor_action}
\end{equation}
with appropriate bundle weights \cite{Mason2005,Movshev2008}. Here ${\cal A}$
is a $(0,1)$ connection, ${\cal F}^{(0,2)}$ is its curvature, and ${\cal B}$
enforces holomorphicity.

The field equations enforce holomorphicity and correspond, under the
Penrose--Ward transform and the usual triviality condition on twistor lines,
to the self-dual gauge equations on spacetime. Equation
\eqref{eq:twistor_action} is not a positive Hermitian curvature norm and does
not alone imply the $L^2$ identity used below.

\subsection{Reconstruction to spacetime and fiber reduction}

Let $A_\mu(x)$ be the reconstructed spacetime gauge field. The declared
reconstruction-and-norm bridge assigns the quadratic functional
\begin{equation}
S_{4}[A]
=
\frac{1}{4g_{\mathrm{eff}}^2}
\int_{Y^4}
F_{\mu\nu}F^{\mu\nu}\,d\mathrm{vol}_4,
\label{eq:4d_action}
\end{equation}
for some effective coupling $g_{\mathrm{eff}}$ determined by the twistor normalization.

\begin{proposition}[Conditional fiber identity]
Assume the reconstruction-and-norm bridge and normalize its chosen Hermitian
fiber measure by \eqref{eq:FSnorm}. Then
the effective coupling satisfies
\begin{equation}
\frac{1}{g_{\mathrm{eff}}^2}
=
\frac{1}{g_{\mathrm{tw}}^2}
\int_{L_x}
\|\psi^{(0)}\|^2\, d\mu_{L_x},
\label{eq:fiber_identity}
\end{equation}
where $\psi^{(0)}$ is the massless twistor harmonic (the fiber representative of the
zero-mode) and $d\mu_{L_x}$ is the measure induced by $\omega_{\mathrm{FS}}$.
\end{proposition}

\begin{proof}
This is the coupling clause of the declared bridge, specialized to the
selected massless representative and the displayed fiber convention.
\end{proof}

\begin{remark}
A derivation rather than an assumption would require a specified twistor
action for the intended spacetime theory, reality structure, Hermitian metric,
gauge fixing, functional measure, and pushforward calculation. None is supplied
by the Fubini--Study area convention alone.
\end{remark}

\subsection{Identification with MTT overlap integrals}

In MTT, the gauge overlap formulation at the matching scale is
\begin{equation}
\frac{1}{g_a^2} = \frac{1}{g_{10}^2} I_a,
\qquad
I_a = \int_{B_a}\langle\omega_a,\omega_a\rangle\, d\mu_{B_a}.
\label{eq:MTT_overlap}
\end{equation}

To compare with \eqref{eq:fiber_identity}, note that in the high-coherence corner the
massless gauge mode factorizes into a spacetime field and a fiber harmonic. Under this
identification,
\[
\int_{L_x}\|\psi^{(0)}\|^2 d\mu_{L_x}
\quad \leftrightarrow \quad
\int_{B_a}\langle\omega_a^{(0)},\omega_a^{(0)}\rangle d\mu_{B_a}
= I_a^{(0)}.
\]

Thus Route~B yields the same functional form as Route~A after imposing the
bridge between the twistor fiber norm and the internal overlap. The
Fubini--Study convention fixes a fiber measure, but does not alone determine
the relative coupling normalization or select the internal metric scale.

\subsection{Explicit SU(2) normalization: recovery of the
\texorpdfstring{$4\pi$}{4 pi} lens factor}

We now carry out the only place in Route~A where a constant could have been ambiguous:
the lens coefficient $\kappa_\ell$.

In Route~A (Paper~II), one used the coefficient $2/(f_2R_{\mathrm{lens}})^2$ in the lens
spectral bound to motivate the round sphere model, yielding
\[
I_2^{(0)} = \mathrm{Area}(S^2_{f_2R_{\mathrm{lens}}}) = 4\pi (f_2R_{\mathrm{lens}})^2.
\]

Route~B reproduces the same factor without reference to that spectral heuristic:

\begin{proposition}[Conditional recovery of $\kappa_\ell=4\pi$]
Let the $SU(2)$ massless gauge mode be constant on the twistor fiber. In
addition to \eqref{eq:FSnorm}, assume the direction-sphere bridge
$dA_{\mathrm{dir}}=2\omega_{\mathrm{FS}}$ and scale that metric by
$(f_2R_{\mathrm{lens}})^2$. Then the induced MTT overlap is
\[
I_2^{(0)} = 4\pi(f_2R_{\mathrm{lens}})^2.
\]
\end{proposition}

\begin{proof}
By the bridge assumption,
$\int_{L_x}dA_{\mathrm{dir}}=2\int_{L_x}\omega_{\mathrm{FS}}=4\pi$.
Metric scaling multiplies this area by $(f_2R_{\mathrm{lens}})^2$,
which gives the stated overlap. The factor of two is part of the explicit
bridge and is not derived from \eqref{eq:FSnorm} alone.
\end{proof}

This recovers the Route~A coefficient conditionally; it does not select the
direction-sphere bridge or lens scale entirely within Route~B.

\subsection{SU(3) remark and Route~B completion criterion}

Twistor theory naturally encodes massless SDYM sectors irrespective of the internal
realization of color, but the explicit reduction of the SU(3) overlap requires specifying
the internal color fiber and its harmonic representative.
Given such a choice (e.g. $\Gamma\backslash\mathrm{Nil}_3$ with left-invariant harmonic 1-forms),
the same bridge \eqref{eq:fiber_identity} applies after the relative
coupling normalization is declared. In particular, the period-normalized
$dx$ or $dy$ representative of Paper~II has squared norm $c$ when $a=b$,
and then Route~B matches Route~A. Imposing unit $L^2$ norm instead would
erase that coefficient unless the bridge weight were changed explicitly.

\begin{remark}[What remains for a fully independent SU(3) Route~B]
To make SU(3) completely independent of Route~A, one must (i) specify the twistor-corner
representation of the color fiber as a canonical twistor bundle, and (ii) compute the
corresponding $L^2$ harmonic norm directly on that bundle. Until those steps are supplied, the $SU(3)$ result is a conditional
cross-check and not an independent normalization theorem.
\end{remark}

\section{Selected-source results and this paper's boundary}
\label{sec:theta-source-consumers}

The eta9 imports are owned contextually by \emph{Flux Compactifications in
Heterotic String Theory}, in the sections on the selected source and the
integral comparison, and completed local tests and the global endpoint.
Gate~1 closes the original-Jacobian campaign at all 30 groups and 225 selected
columns \cite{ThetaGate1}. This is an exact finite source calculation, not
the global integral meridian, the 248-coordinate readout or the analytic
Deligne class $\beta_{\mathbb C}$.

The completed three-cycle transport has independent cycles with Gram matrix
$-2I_3$ and zero detected affine pairings \cite{ThetaCycles}. Its scope is
a non-detecting subsystem: it does not prove global triviality of the twist,
and rerunning that same subsystem is not the missing detection theorem.
Neither this conclusion nor the finite gate identifies an auxiliary
Circle--Lens--Nil model with the physical q79 topology.

The large-gauge result is owned by \emph{Cohesive Closure Repair and Its
Hodge, Kernel, and Projection Shadows}, in its subsection on the rank-zero
large-gauge kernel. It gives an injective immersion
$b_{K3}:\mathbb R^2\to\mathbb T^{20}$ with dense nonclosed image
\cite{ThetaBK3}. Rank zero counts periodic gauge identifications, not source
coordinates or physical modes. It removes a kernel ambiguity without
requiring numerical periods, but does not provide kinetic normalization,
the selected shared-circle action or the unrelated eta9 affine lift.

For Route~B, neither the finite quotient closure nor the global gauge-kernel
theorem selects a Hermitian fiber norm, a direction-sphere bridge, or the
period-normalized color representative. Equality with Route~A remains a
conditional dictionary check. A physical q79 source and coefficient-level
remainder control would still be needed to turn that agreement into an
independent prediction.


The managed evidence block below retains its earlier frozen profile snapshot.
Its historical source-status labels do not supersede the current exact,
local and open boundaries just stated or reopen the retained hidden-HYM
existence and shared-primitive Standard Model results.

% BEGIN MTT MANAGED COMPUTATIONAL EVIDENCE
\section*{Computational Evidence and Reproducibility}
The HYM, finite-matrix, and common-scheme precision rows provide the computational target for the conditional twistor-action audit. They do not supply the missing action normalization or prove that the twistor corner is the physical branch. Other sector rows are context, and the strict upgrade remains open.

The referenced rows are frozen to the curated results repository state identified below. Hashes are grouped in eight-character blocks for line breaking.
\begin{quote}\small
\textbf{Repository:} \url{https://github.com/PeterNero/mtt-results-repro}\\
\textbf{Commit:} \texttt{31247ebb 5c22f3fb b5443024 365433c6 ee0bff4a}\\
\textbf{Manifest:} \path{release/result_manifest.json}\\
\textbf{Manifest SHA-256:}\\
\texttt{fb399689 60b00584 631dbf53 1a708e18}\\
\texttt{ef928d6b 6d935119 c185d7f6 32b1e7cd}
\end{quote}

Tier labels are quoted verbatim from that manifest. A row used directly supports only the specific computational statement identified above; a corpus-state cross-check does not prove this paper's local theorems; and an open row is evidence of an unresolved obligation, never of closure.

\par\begingroup\dimen0=\pagegoal\advance\dimen0 by -\pagetotal\ifdim\dimen0<7\baselineskip\newpage\fi\endgroup
\paragraph{Rows used directly in this paper.}
\begin{itemize}
\setlength{\itemsep}{0.25em}
\item \path{A19/hym_wiener_contraction} (\emph{numeric certified}).\par Weighted-theta Fourier-tail and Wiener contraction certificate.
\item \path{A02/precision_15_source_transport} (\emph{profile replay}).\par Fifteen measured source coordinates, Jacobian and covariance transport.
\item \path{A02/precision_8x8_workspace} (\emph{profile replay}).\par Eight-coordinate SMDR output with positive-definite 8x8 covariance.
\item \path{A01/qutrit_weyl_27_matrix} (\emph{derived exact}).\par Sparse 27x27 qutrit-Weyl left-action realization.
\end{itemize}

\par\begingroup\dimen0=\pagegoal\advance\dimen0 by -\pagetotal\ifdim\dimen0<7\baselineskip\newpage\fi\endgroup
\paragraph{Corpus-state cross-checks.}
\begin{itemize}
\setlength{\itemsep}{0.25em}
\item \path{A01/charged_yukawa_higgs_profile} (\emph{profile replay}).\par Versioned Yu, Yd, Ye and lambda\_H profile packet.
\item \path{A14/ckm_prediction_profile} (\emph{numeric certified}).\par Three selected CKM profile rows and uncertainty comparison.
\item \path{A01/current_global_lock} (\emph{profile replay}).\par Current non-looping global status and source-certificate map.
\item \path{A01/direct_k_higgs_row} (\emph{derived exact}).\par Promoted direct K\_threshold.Omega\_H.lambda row.
\item \path{A22/e6_qpsi_qcd_anomaly} (\emph{derived exact}).\par E6 Qpsi matter/exotic QCD anomaly cancellation audit.
\item \path{A04/final_12_of_12_audit} (\emph{profile replay}).\par Twelve-obligation embedded renormalized-SM equivalence audit.
\item \path{A40/neutral_two_primitive_profile} (\emph{profile replay}).\par Measured two-splitting neutrino profile, masses and Dirac Yukawa rows.
\item \path{A01/strict_pew_row} (\emph{derived exact}).\par Promoted P\_EW source row at the declared one-shared-primitive standard.
\end{itemize}

\par\begingroup\dimen0=\pagegoal\advance\dimen0 by -\pagetotal\ifdim\dimen0<7\baselineskip\newpage\fi\endgroup
\paragraph{Open boundary (not evidence of closure).}
\begin{itemize}
\setlength{\itemsep}{0.25em}
\item \path{A05/strict_upgrade_ledger} (\emph{open}).\par Current 2/9 strict no-knob upgrade ledger.
\end{itemize}

No imported row changes theorem ownership or promotes a neighboring claim: all local statements retain their stated hypotheses, domains, and limitations.
% END MTT MANAGED COMPUTATIONAL EVIDENCE

\begin{thebibliography}{99}

\bibitem{ThetaGate1}
P.~Nero, \emph{Selected eta9 original-Jacobian Gate 1 campaign} (2026).
Frozen source at curated manifest commit
\texttt{f141a20e a23c5c3f f19cc216 1c0e226e 29ade8a7}:
\href{https://github.com/PeterNero/mtt-results-repro/blob/f141a20ea23c5c3ff19cc2161c0e226e29ade8a7/release/results/eta9_gate1_campaign/artifact.json}{source artifact}.

\bibitem{ThetaCycles}
P.~Nero, \emph{Complete selected three-cycle non-detection decision} (2026).
Frozen source at curated manifest commit
\texttt{f141a20e a23c5c3f f19cc216 1c0e226e 29ade8a7}:
\href{https://github.com/PeterNero/mtt-results-repro/blob/f141a20ea23c5c3ff19cc2161c0e226e29ade8a7/release/results/eta9_three_cycle_b96_nondetection/artifact.json}{source artifact}.

\bibitem{ThetaBK3}
P.~Nero, \emph{Selected K3 rank-zero large-gauge kernel} (2026).
Frozen source at curated manifest commit
\texttt{f141a20e a23c5c3f f19cc216 1c0e226e 29ade8a7}:
\href{https://github.com/PeterNero/mtt-results-repro/blob/f141a20ea23c5c3ff19cc2161c0e226e29ade8a7/release/results/q79_bk3_rank_zero_kernel/artifact.json}{source artifact}.

\bibitem{Mason2005}
L.~J. Mason,
\emph{Twistor actions for non-self-dual fields: a derivation of
twistor-string theory},
JHEP \textbf{10} (2005) 009.
\url{https://doi.org/10.1088/1126-6708/2005/10/009}

\bibitem{Movshev2008}
M.~V. Movshev,
\emph{A note on self-dual Yang--Mills theory},
\url{https://arxiv.org/abs/0812.0224}

\bibitem{MTT-Admissibility}
P.~Nero,
\emph{Modal Triplet Theory: Admissibility, Encodings, and the Structure of Physical Description},
Zenodo preprint, January 2026.
\url{https://doi.org/10.5281/zenodo.18255621}

\bibitem{MTT-Foundation}
P.~Nero,
\emph{Modal Triplet Theory: Foundation},
Zenodo preprint, September 2025.
\url{https://doi.org/10.5281/zenodo.16949762}

\bibitem{MTT-FixedPoints}
P.~Nero,
\emph{Fixed Points I--VI: Complete Coherence Spine},
Zenodo preprints, August 2025.
\url{https://doi.org/10.5281/zenodo.16948748}

\bibitem{MTT-ProjectionAdmissibility}
P.~Nero,
\emph{The Projection--Admissibility Principle: Structural Constraints on Effective Physical Description},
Zenodo preprint, January 2026.
\url{https://doi.org/10.5281/zenodo.18255838}

\bibitem{MTT-Closure}
P.~Nero,
\emph{Closure and Inevitability in Modal Triplet Theory},
Zenodo preprint, January 2026.
\url{https://doi.org/10.5281/zenodo.18255510}

\bibitem{MTT-CoherenceCapacity}
P.~Nero,
\emph{Coherence Capacity as the Fundamental Resource of Effective Physics},
Zenodo preprint, January 2026.
\url{https://doi.org/10.5281/zenodo.18255905}

\bibitem{MTT-CoherenceDynamics}
P.~Nero,
\emph{Dynamics of Coherence Capacity: Transport, Concentration, and Exhaustion},
Zenodo preprint, January 2026.
\url{https://doi.org/10.5281/zenodo.18256048}

\bibitem{MTT-QM}
P.~Nero,
\emph{Modal Triplet Theory: From MTT to Quantum Mechanics},
Zenodo preprint, September 2025.
\url{https://doi.org/10.5281/zenodo.17074246}

\bibitem{MTT-QFT}
P.~Nero,
\emph{From MTT to Quantum Field Theory},
Zenodo preprint, 2025.
\url{https://doi.org/10.5281/zenodo.17068816}

\bibitem{MTT-GR}
P.~Nero,
\emph{Modal Triplet Theory: From MTT to General Relativity},
Zenodo preprint, October 2025.
\url{https://doi.org/10.5281/zenodo.16950597}

\bibitem{MTT-UVQG}
P.~Nero,
\emph{Modal Triplet Theory: From MTT to a UV-Finite, Unitary Quantum Gravity},
Zenodo preprint, 2025.
\url{https://doi.org/10.5281/zenodo.17077671}

\bibitem{MTT-Measurement}
P.~Nero,
\emph{Measurement as Disturbance and Stabilization in Modal Triplet Theory},
Zenodo preprint, 2025.
\url{https://doi.org/10.5281/zenodo.17177404}

\bibitem{MTT-Irreversibility}
P.~Nero,
\emph{Projection, Probability, and Irreversibility: Shadow Bridges Between Measurement, Black Holes, and Cosmology in Modal Triplet Theory},
Zenodo preprint, January 2026.
\url{https://doi.org/10.5281/zenodo.18256408}

\bibitem{MTT-Bell-Beables}
P.~Nero,
\emph{Modal Fixed Points, Bell’s Beables, and the Limits of Factorization},
Zenodo preprint, 2025.
\url{https://doi.org/10.5281/zenodo.17076300}

\bibitem{MTT-Bell-Temporal}
P.~Nero,
\emph{Temporal Bell Inequalities and Global Consistency in Modal Triplet Theory},
Zenodo preprint, August 2025.
\url{https://doi.org/10.5281/zenodo.18208884}

\bibitem{MTT-Stochastic}
P.~Nero,
\emph{From Modal Triplet Theory to Indivisible Stochastic Processes: A First-Principles, Fully Rigorous Derivation},
Zenodo preprint, January 2026.
\url{https://doi.org/10.5281/zenodo.18254862}

\end{thebibliography}
\end{document}
